\documentclass[11pt]{article}

\usepackage{amsmath}
\usepackage{amssymb}
\usepackage{cite}
\usepackage{hyperref}
\usepackage{lineno}

% Add bibiography style in the brackets below.
%\bibliographystyle{}

% The title text is entered in the brackets below
\title{Article title}

% Author names are listed here.
\author{
 Author, First\\
  \texttt{first1.last1@xxxxx.com}
  \and
  Author, Second\\
  \texttt{first2.last2@xxxxx.com}
}
\date{}

\begin{document}
\maketitle



\section*{Abstract}
Abstract text.

\section*{Introduction}

Introduction text.


\section*{Materials and Methods}

Materials and methods text.

This is what an equation might look like.
\begin{equation}
x = \pm \sqrt{\frac{b^2 - 4ac}{2a}};
\end{equation}
\section*{Results}

Results text.


\subsection*{Example of a Subsection}


\section*{Discussion}



\section*{Acknowledgments}


\section*{References}
% If you did not use a .bib file, but instead included your references
% in here, below is some example code:
%\begin{thebibliography}{6}

%\bibitem{mascot} Jack and Jill, {\em A Journal}, The Big Journal,{\bf 22}, 332 -- 333
% \end{thebibliography}


% Should you use a .bib file, below is some sample code:
%\bibliographystyle{plain}
%\bibliography{sample1,sample2,...,samplen}


\section*{Figure Legends}

\begin{figure}
% Below is the command to include your figure file.
%\includegraphics[width=0.5\textwidth]{your_figure_name}}
\caption{Figure caption.}
\label{Figure_label}
\end{figure}


\section*{Tables}


\begin{table}[h!]
  \begin{center}
    \begin{tabular}{| l c r |}
    \hline
    1 & 2 & 3 \\
    4 & 5 & 6 \\
    7 & 8 & 9 \\
    \hline
    \end{tabular}
  \end{center}
  \caption{A simple table}
\end{table}



\end{document}

